\documentclass[10pt,a4paper]{article}
\usepackage[T1]{fontenc}
\usepackage[utf8]{inputenc}
\usepackage[margin=2cm,landscape]{geometry}
\usepackage{amsmath,amssymb}
\usepackage{array,booktabs,longtable}
\usepackage{xcolor}
\usepackage{enumitem}
\usepackage[hidelinks]{hyperref}
\newcommand{\lean}[1]{\textcolor{blue!50!black}{\texttt{\detokenize{#1}}}}
\setlength{\LTpre}{6pt}\setlength{\LTpost}{6pt}
\renewcommand{\arraystretch}{1.25}
\begin{document}


\begin{center}{\Large\bfseries Thin perfect matchings [HOG26]: Lean formalisation}\\[4pt]

{\small Correspondence between the paper and the Lean~4/Mathlib formalisation}\end{center}

\noindent\textbf{Paper.} A.~Haqi, S.~Oveis Gharan. \emph{On Thin Perfect Matchings up to Polylogarithmic Factors}. arXiv:2606.01330v2 (2026) [HOG26].\par\smallskip

The library \texttt{ThinMatching/} (independent of \texttt{GsimOT}) formalises [HOG26]: Lemma~1.2, Theorem~1.3 (support-constrained) and Theorem~1.4 (relaxed, bipartite and general) with all lemmas of Sections~3 and~4. Theorem~1.4 is the cited Theorem~35 of the main paper (p.~32). A proof overview is in \texttt{ThinMatching/HOG26\_overview.pdf}.\par\paragraph{Checking.} Run \texttt{bash build.sh} in the repository root (see \texttt{README.md}). The script builds the library and runs \texttt{scripts/CheckAxioms.lean}, which must report only the standard axioms \texttt{propext}, \texttt{Classical.choice}, \texttt{Quot.sound}: no \texttt{sorry} and no additional axioms.

\paragraph{Main statement in Lean.} \lean{ThinMatching.theorem14_bip}: assuming [RS14] (\lean{RackeShahTheorem}), there is an absolute $C>0$ such that
for every bipartite fractional perfect matching $x$ on $L\cup R$, $|L|=|R|=n$, some perfect matching $M\subseteq L\times R$
satisfies $|M\cap\delta(S)|\le C\ln(2n)\ln\ln(4n)\,x(\delta(S))$ for all $S\subseteq L\cup R$.
\lean{ThinMatching.theorem13_support}: assuming Lemma~3.6 (\lean{Lemma36}), every $x\in\mathrm{PM}(G)$ has a
$C|V|\ln(2|V|)$-thin perfect matching inside $\mathrm{supp}(x)$.

\section*{Correspondence}

\begin{longtable}{>{\raggedright\arraybackslash}p{3.2cm} >{\centering\arraybackslash}p{1.1cm} >{\raggedright\arraybackslash}p{8.4cm} >{\raggedright\arraybackslash}p{10.2cm}}
\toprule
\textbf{Paper} & \textbf{Page} & \textbf{Statement} & \textbf{Lean declaration(s) and location} \\
\midrule\endhead
\bottomrule\endfoot
{}Sec.~1, Sec.~2 & 1, 4 & Weighted graph $G=(V,E,x)$, $x(\delta(S))$, perfect matchings, $|M\cap\delta(S)|$, $\alpha$-thin, $M\subseteq\mathrm{supp}(x)$ & \lean{ThinMatching.WGraph}\newline{\scriptsize\texttt{ThinMatching/Basic.lean:30}} \newline \lean{ThinMatching.cutWeight}\newline{\scriptsize\texttt{ThinMatching/Basic.lean:41}} \newline \lean{ThinMatching.PerfectMatching}\newline{\scriptsize\texttt{ThinMatching/Basic.lean:44}} \newline \lean{ThinMatching.crossCount}\newline{\scriptsize\texttt{ThinMatching/Basic.lean:50}} \newline \lean{ThinMatching.IsThin}\newline{\scriptsize\texttt{ThinMatching/Basic.lean:53}} \newline \lean{ThinMatching.InSupport}\newline{\scriptsize\texttt{ThinMatching/Basic.lean:57}} \\ \midrule
{}Sec.~1, Thm.~2.5 & 1, 5 & $x\in\mathrm{PM}(G)$: convex combination of perfect matchings; bipartite version (degree constraints) & \lean{ThinMatching.IsFracPM}\newline{\scriptsize\texttt{ThinMatching/Basic.lean:60}} \newline \lean{ThinMatching.IsFracPMBip}\newline{\scriptsize\texttt{ThinMatching/Basic.lean:84}} \newline \lean{ThinMatching.bipGraph}\newline{\scriptsize\texttt{ThinMatching/Basic.lean:90}} \\ \midrule
{}Thm.~1.4 (bipartite) & 1, 2 & Bipartite cut weight, crossing count and thinness for $M=\{(l,\pi l)\}$ & \lean{ThinMatching.bipCutWeight}\newline{\scriptsize\texttt{ThinMatching/Basic.lean:71}} \newline \lean{ThinMatching.bipCrossCount}\newline{\scriptsize\texttt{ThinMatching/Basic.lean:75}} \newline \lean{ThinMatching.IsThinBip}\newline{\scriptsize\texttt{ThinMatching/Basic.lean:79}} \\ \midrule
{}Lemma~1.2 & 1 & Every perfect matching in $\mathrm{supp}(x)$ of the graph of Fig.~1 ($2n+2$ vertices) is at best $\frac n2$-thin & \lean{ThinMatching.lemma12}\newline{\scriptsize\texttt{ThinMatching/LowerBound.lean:200}} \newline \lean{ThinMatching.lemma12_not_thin}\newline{\scriptsize\texttt{ThinMatching/LowerBound.lean:208}} \\ \midrule
{}Fig.~1 & 1 & The graph: solid edges $1-\frac1n$, dotted edges $\frac1n$; $x$ as the average of $n$ perfect matchings & \lean{ThinMatching.LowerBound.fig1}\newline{\scriptsize\texttt{ThinMatching/LowerBound.lean:43}} \newline \lean{ThinMatching.LowerBound.Mk}\newline{\scriptsize\texttt{ThinMatching/LowerBound.lean:74}} \newline \lean{ThinMatching.LowerBound.fig1_isFracPM}\newline{\scriptsize\texttt{ThinMatching/LowerBound.lean:83}} \newline \lean{ThinMatching.LowerBound.fig1_bad_cut}\newline{\scriptsize\texttt{ThinMatching/LowerBound.lean:135}} \\ \midrule
{}Theorem~1.3 & 2 & Support-constrained $O(n\log n)$-thin perfect matching (assuming Lemma~3.6) & \lean{ThinMatching.theorem13_support}\newline{\scriptsize\texttt{ThinMatching/SupportConstrained.lean:212}} \\ \midrule
{}Def.~3.1, Lemma~3.2 & 5 & Trimming $\tilde x$; removed weight $\le(|V|-1)\varepsilon$; components $\varepsilon$-edge-connected & \lean{ThinMatching.trimGraph}\newline{\scriptsize\texttt{ThinMatching/SupportConstrained/Trimming.lean:45}} \newline \lean{ThinMatching.IsTrimmed}\newline{\scriptsize\texttt{ThinMatching/SupportConstrained/Trimming.lean:52}} \newline \lean{ThinMatching.exists_split}\newline{\scriptsize\texttt{ThinMatching/SupportConstrained/Trimming.lean:84}} \newline \lean{ThinMatching.trimming_exists}\newline{\scriptsize\texttt{ThinMatching/SupportConstrained/Trimming.lean:156}} \newline \lean{ThinMatching.lemma32}\newline{\scriptsize\texttt{ThinMatching/SupportConstrained/Trimming.lean:213}} \\ \midrule
{}Algorithm~1 & 6 & Independent sampling of $\alpha$ copies of each edge with probability $\tilde x_e$ (finite product space) & \lean{ThinMatching.bernW}\newline{\scriptsize\texttt{ThinMatching/SupportConstrained/Prob.lean:27}} \newline \lean{ThinMatching.prob}\newline{\scriptsize\texttt{ThinMatching/SupportConstrained/Prob.lean:31}} \newline \lean{ThinMatching.sampleProb}\newline{\scriptsize\texttt{ThinMatching/SupportConstrained/Sampling.lean:67}} \newline \lean{ThinMatching.sampleCut}\newline{\scriptsize\texttt{ThinMatching/SupportConstrained/Sampling.lean:81}} \newline \lean{ThinMatching.sampleGraph}\newline{\scriptsize\texttt{ThinMatching/SupportConstrained/Sampling.lean:85}} \\ \midrule
{}Thm.~2.4, Def.~3.3 & 5 & Tutte's theorem (Mathlib) and Tutte witnesses & \lean{ThinMatching.IsTutteWitness}\newline{\scriptsize\texttt{ThinMatching/SupportConstrained/TutteWitness.lean:50}} \newline \lean{ThinMatching.witnessF}\newline{\scriptsize\texttt{ThinMatching/SupportConstrained/TutteWitness.lean:62}} \newline \lean{ThinMatching.exists_witness_of_no_pm}\newline{\scriptsize\texttt{ThinMatching/SupportConstrained/TutteWitness.lean:170}} \\ \midrule
{}Lemma~3.4 & 6 & At most $(|V|+1)^{|V|}$ Tutte witnesses & \lean{ThinMatching.card_witnesses_le}\newline{\scriptsize\texttt{ThinMatching/SupportConstrained/TutteWitness.lean:86}} \\ \midrule
{}Lemma~3.5, Eq.~(1) & 6 & The sample contains a perfect matching w.h.p.; $x(F)\ge1$ (odd-set argument) & \lean{ThinMatching.lemma35}\newline{\scriptsize\texttt{ThinMatching/SupportConstrained/Lemma35.lean:62}} \newline \lean{ThinMatching.one_le_witnessF_weight}\newline{\scriptsize\texttt{ThinMatching/SupportConstrained/TutteWitness.lean:134}} \newline \lean{ThinMatching.witness_bound_le}\newline{\scriptsize\texttt{ThinMatching/SupportConstrained/Lemma35.lean:122}} \\ \midrule
{}Lemma~3.6 & 7 & Sampling preserves cuts up to $O(\alpha)$ w.h.p.\ (proof omitted in [HOG26]) --- \textbf{explicit hypothesis} & \lean{ThinMatching.Lemma36}\newline{\scriptsize\texttt{ThinMatching/SupportConstrained/Sampling.lean:120}} \\ \midrule
{}Proof of Thm.~1.3 & 7 & Union over components; thinness for all cuts; $M\subseteq Y\subseteq\mathrm{supp}(\tilde x)\subseteq\mathrm{supp}(x)$ & \lean{ThinMatching.sum_cut_eq_sum_blocks}\newline{\scriptsize\texttt{ThinMatching/SupportConstrained.lean:72}} \newline \lean{ThinMatching.crossCount_le_of_good}\newline{\scriptsize\texttt{ThinMatching/SupportConstrained.lean:119}} \newline \lean{ThinMatching.prob_bad_blocks_le}\newline{\scriptsize\texttt{ThinMatching/SupportConstrained.lean:174}} \\ \midrule
{}Def.~2.1 & 4 & Weighted tree with Steiner nodes (rooted, parent function); $y(\delta_{\mathcal T}(\mathcal S))$; $\lambda_{\mathcal T}(S)$; $q$-tree-cut-sparsifier & \lean{ThinMatching.RootedTree}\newline{\scriptsize\texttt{ThinMatching/TreeSparsifier.lean:29}} \newline \lean{ThinMatching.subtree}\newline{\scriptsize\texttt{ThinMatching/TreeSparsifier.lean:39}} \newline \lean{ThinMatching.treeCutWeight}\newline{\scriptsize\texttt{ThinMatching/TreeSparsifier.lean:44}} \newline \lean{ThinMatching.lambdaT}\newline{\scriptsize\texttt{ThinMatching/TreeSparsifier.lean:50}} \newline \lean{ThinMatching.IsTreeCutSparsifier}\newline{\scriptsize\texttt{ThinMatching/TreeSparsifier.lean:55}} \\ \midrule
{}Thm.~2.2 [RS14] & 4 & Every graph has an $O(\log n\log\log n)$-tree-cut-sparsifier --- \textbf{explicit hypothesis} & \lean{ThinMatching.RackeShahTheorem}\newline{\scriptsize\texttt{ThinMatching/TreeSparsifier.lean:63}} \\ \midrule
{}Defs.~4.1, 4.2, 4.4 & 7--8 & Tree cuts, load of a matching on a tree edge, discrepancy $\Delta$ & \lean{ThinMatching.load}\newline{\scriptsize\texttt{ThinMatching/Load.lean:32}} \newline \lean{ThinMatching.loadBip}\newline{\scriptsize\texttt{ThinMatching/Load.lean:46}} \newline \lean{ThinMatching.discBip}\newline{\scriptsize\texttt{ThinMatching/Load.lean:50}} \newline \lean{ThinMatching.discNonBip}\newline{\scriptsize\texttt{ThinMatching/Load.lean:36}} \\ \midrule
{}Proposition~4.3 & 7 & $\mathrm{load}_M(f)\le\beta\,x(\delta(V_r(f)))$ for all tree edges $\Rightarrow$ $M$ is $\beta q$-thin & \lean{ThinMatching.prop43}\newline{\scriptsize\texttt{ThinMatching/LoadVsCut.lean:170}} \newline \lean{ThinMatching.LoadVsCut.exists_cut_edge}\newline{\scriptsize\texttt{ThinMatching/LoadVsCut.lean:107}} \newline \lean{ThinMatching.LoadVsCut.treeCutWeight_subtree}\newline{\scriptsize\texttt{ThinMatching/LoadVsCut.lean:148}} \\ \midrule
{}Lemma~4.5 & 8 & $x(\delta(S))\ge\Delta_{\mathrm{bip}}(S)$; $x(\delta(S))\ge1$ for $|S|$ odd & \lean{ThinMatching.lemma45_bip}\newline{\scriptsize\texttt{ThinMatching/Discrepancy.lean:21}} \newline \lean{ThinMatching.lemma45}\newline{\scriptsize\texttt{ThinMatching/Discrepancy.lean:81}} \newline \lean{ThinMatching.one_le_crossCount_of_odd}\newline{\scriptsize\texttt{ThinMatching/Discrepancy.lean:43}} \newline \lean{ThinMatching.cutWeight_eq_sum_crossCount}\newline{\scriptsize\texttt{ThinMatching/Discrepancy.lean:58}} \\ \midrule
{}Lemma~4.6, Algorithm~2 & 9--10 & Tree pairing: a perfect matching with $\mathrm{load}_M(f)=\Delta(V_r(f))$ for every tree edge & \lean{ThinMatching.lemma46_bip}\newline{\scriptsize\texttt{ThinMatching/TreePairing.lean:135}} \newline \lean{ThinMatching.lemma46}\newline{\scriptsize\texttt{ThinMatching/TreePairing.lean:212}} \newline \lean{ThinMatching.subtree_laminar}\newline{\scriptsize\texttt{ThinMatching/TreePairing.lean:56}} \newline \lean{ThinMatching.subtree_ssubset_of_separates}\newline{\scriptsize\texttt{ThinMatching/TreePairing.lean:70}} \\ \midrule
{}Theorem~1.4 (bipartite) & 2 & An $O(\log n\log\log n)$-thin perfect matching $M\subseteq L\times R$ (assuming [RS14]) & \lean{ThinMatching.theorem14_bip}\newline{\scriptsize\texttt{ThinMatching/Theorem14.lean:135}} \newline \lean{ThinMatching.Theorem14.bipPM}\newline{\scriptsize\texttt{ThinMatching/Theorem14.lean:30}} \newline \lean{ThinMatching.Theorem14.cutWeight_bipGraph}\newline{\scriptsize\texttt{ThinMatching/Theorem14.lean:64}} \newline \lean{ThinMatching.Theorem14.log_mul_loglog_le}\newline{\scriptsize\texttt{ThinMatching/Theorem14.lean:96}} \\ \midrule
{}Theorem~1.4 (general) & 2 & An $O(\log n\log\log n)$-thin perfect matching $M\subseteq\binom V2$ for $x\in\mathrm{PM}(G)$ (assuming [RS14]) & \lean{ThinMatching.theorem14}\newline{\scriptsize\texttt{ThinMatching/Theorem14.lean:180}} \\ \midrule
{}Main paper, Thm.~35, Eq.~(23) & -- & Bridge to \texttt{GsimOT}: Theorem~35 (p.~32, thinness Eq.~(23)) and Theorem~11 of the main paper assuming only [RS14] & \lean{GsimOT.thinPermutationTheorem_of_rackeShah}\newline{\scriptsize\texttt{Theorem13\_CutUpperBound/FromThinMatching.lean:18}} \newline \lean{GsimOT.theorem13_of_rackeShah}\newline{\scriptsize\texttt{Theorem13\_CutUpperBound/FromThinMatching.lean:28}} \newline \lean{GsimOT.theorem13_power_of_rackeShah}\newline{\scriptsize\texttt{Theorem13\_CutUpperBound/FromThinMatching.lean:38}} \\
\end{longtable}

\section*{Notes on the formalisation}

\begin{enumerate}[leftmargin=*]
\item \textbf{Independence.} The library \texttt{ThinMatching} imports only Mathlib and does not depend on the \texttt{GsimOT} formalisation. It has its own definitions of weighted graphs, cuts, perfect matchings, thinness and trees. The only connection is the bridge file \texttt{GsimOT/Theorem13\_CutUpperBound/FromThinMatching.lean}, which derives the main paper's Theorem~35 from the bipartite case of Theorem~1.4. The main paper's $\alpha$-thinness, Eq.~(23) (cuts $(X\cup Y,\overline X\cup\overline Y)$ of $K_{n,n}$, all $X,Y\subseteq[n]$), is exactly \lean{IsThinBip} with $S_L=X$, $S_R=Y$, and its constant $C_{\mathrm{thin}}\ln(2n)\ln\ln(4n)$ is the one of \lean{theorem14_bip} (same $C$).
\item \textbf{Cited results.} Theorem~2.2 (R\"acke--Shah [RS14], tree-cut sparsifiers) is used but not proved in [HOG26]. It is the explicit hypothesis \lean{RackeShahTheorem}, with the bound $C\ln(2n)\ln\ln(4n)$, which is positive for all $n\ge1$. Lemma~3.6 is stated in [HOG26] with its proof omitted (``a simple Chernoff bound and a cut-counting argument''). It is the explicit hypothesis \lean{Lemma36}, formulated in the finite sampling model of Algorithm~1, for any $\alpha\ge cN\ln(2N)$ with $N\ge|W|$, with failure probability $\le1/(4N)$. This is what the omitted argument (Karger's cut counting + Chernoff) yields. No Lean \texttt{axiom} is introduced; Theorems~1.3 and 1.4 take these hypotheses as arguments.
\item \textbf{Trees.} A weighted tree with Steiner nodes is represented as a rooted tree on a finite type via a parent function, with $V\hookrightarrow\mathcal V$; the tree edges are $\{v,\mathrm{parent}(v)\}$. Every tree can be rooted, so this is no loss of generality. The tree path between two vertices uses the edge $\{v,\mathrm{parent}(v)\}$ iff exactly one of them lies in the subtree of $v$; this is how loads are defined. Discrepancies are taken on the subtree side $V_s(f)$, which equals the root side $V_r(f)$ of the paper since $|L|=|R|$ resp.\ $|V|$ is even.
\item \textbf{Lemma 4.6 (Algorithm 2).} Instead of executing the greedy vertex by vertex, the proof matches, one pair at a time, a matchable pair inside a \emph{minimal} subtree that contains one. This is the first pair the greedy forms. The paper's main observation (no proper subtree of it contains a matchable pair) shows that load and discrepancy change in the same way, which gives $\mathrm{load}=\Delta$ by induction.
\item \textbf{Proposition 4.3.} The charging step uses a lemma (\lean{exists_cut_edge}) finding, for $a\in\mathcal S$, $b\notin\mathcal S$, a tree edge cut by $\mathcal S$ that separates $a$ and $b$, by induction on the depths instead of via the explicit path. Lemma~4.5 in the general case is derived from the convex combination of perfect matchings (each crosses an odd set), the easy direction of Theorem~2.5.
\item \textbf{Theorem 1.3.} The paper argues per component. The formalisation also checks cuts through several components: sampled edges lie inside blocks, so the per-component bounds add up. Further choices: a union bound over at most $|V|$ components, $\varepsilon=1/(4|V|)$ (the paper's $1/(8n)$ for $|V|=2n$), and $x(F)\ge1$ (slightly stronger than the paper's route). Trimming is modelled by block labellings, and each step adds one block.
\item \textbf{Not formalised.} Corollary~1.6 (metric distortion) depends on [ACR23, Prop.~3] and on mechanism definitions outside the mathematical core of the paper. Theorem~2.3 (algorithmic sparsifiers) is not used.
\item Lemma~1.2 is proved for every $n\ge1$; as in the paper, one crossing matching edge is counted (ratio $\ge n/2$).
\end{enumerate}

\end{document}